\documentclass[../NDCNNAI_main.tex]{subfiles}
\graphicspath{{\subfix{../fig/}}}

\begin{document}

\subsection{Введение и предварительные сведения}

В данном разделе мы обрисуем общую постановку задачи обучения с учителем, введём остаточные сети (ResNet) и их интерпретацию как динамических систем, а также обсудим ключевые понятия \emph{пространства гипотез}, \emph{потокового отображения} и различия между \emph{точным представлением} и \emph{аппроксимацией}.

Рассматривается классическая задача обучения с учителем \cite{Li2023DeepLD}. Пусть заданы:
\begin{itemize}
    \item Входное пространство $\mathcal{X} \subset \mathbb{R}^n$.
    \item Выходное пространство $\mathcal{Y} \subset \mathbb{R}^m$.
    \item Неизвестная целевая функция $F: \mathcal{X} \to \mathcal{Y}$.
    \item Набор данных $\{(x_i, y_i)\}_{i=1}^{N}$, где $y_i = F(x_i)$.
\end{itemize}
Цель --- найти функцию $G$ из некоторого \emph{пространства гипотез} $\mathcal{H}$, которая хорошо приближает $F$ на данных. Это приводит к задаче минимизации \emph{эмпирического риска}:
\begin{equation*}
\inf_{G \in \mathcal{H}} \frac{1}{N} \sum_{i=1}^{N} \ell(F(x_i), G(x_i)),
\end{equation*}
где $\ell$ --- функция потерь (например, квадратичная или кросс-энтропия) \cite[Раздел 1.1]{Li2023DeepLD}.

В данном подразделе пространство гипотез $\mathcal{H}$ будет строиться на основе \emph{потоковых отображений} обыкновенных дифференциальных уравнений (ОДУ), что является идеализацией глубоких остаточных сетей.

\textbf{Дискретная остаточная сеть (ResNet)} \cite[Раздел 1.2]{Li2023DeepLD} задаётся разностным уравнением:
\begin{equation}
\label{eq:resnet_discrete}
z_{s+1} = z_s + \delta f_{\theta_s}(z_s), \quad s = 0, \ldots, S-1, \quad z_0 = x,
\end{equation}
где $\delta > 0$ --- размер шага (часто $\delta=1$), $f_{\theta_s}$ --- нелинейное преобразование соответствующего слоя (например, аффинное отображение с последующей функцией активации), а $S$ --- общая глубина сети.

\textbf{Непрерывный предел.} Интерпретируя индекс слоя $s$ как дискретное время, можно рассмотреть непрерывный предел при $\delta \to 0$, полагая $t = s\delta$, $T = S\delta$. Тогда \eqref{eq:resnet_discrete} переходит в ОДУ \cite[Раздел 1.3]{Li2023DeepLD}:
\begin{equation}
\label{eq:resnet_ode}
\dot{z}(t) = f_{\theta(t)}(z(t)), \quad t \in [0, T], \quad z(0) = x.
\end{equation}
Если $f_{\theta(t)}$ глобально липшицева по $z$, то решение \eqref{eq:resnet_ode} существует и единственно. \emph{Потоковое отображение} (flow map) в момент времени $T$ определяется как $\varphi_T(\cdot; \theta): x \mapsto z(T)$.

Таким образом, глубокую ResNet можно рассматривать как дискретизацию по времени некоторой управляемой динамической системы. Этот подход позволяет применять для анализа глубоких сетей достаточной развитый аппарат теории ОДУ, оптимального управления и дифференциальных уравнений.

Исходя из двух представлений (дискретного и непрерывного), вводятся соответствующие пространства гипотез.

\begin{definition}[Пространство гипотез дискретных ResNet]
Для заданных глубины $S$ и шага $\delta$ определим
\begin{equation*}
\mathcal{H}_{\text{resnet}}(S,\delta) = \{ g \circ \varphi_{S,\delta}(\cdot; \theta) : g \in \mathcal{G}, \theta = (\theta_0, \ldots, \theta_{S-1}) \},
\end{equation*}
где $\varphi_{S,\delta}$ --- потоковое отображение дискретной системы \eqref{eq:resnet_discrete}, а $\mathcal{G}$ --- \emph{терминальное семейство} (обычно простые функции, например, аффинные отображения $\mathbb{R}^n \to \mathcal{Y}$).
\end{definition}

\begin{definition}[Пространство гипотез непрерывных ODE-сетей]
Для заданного горизонта времени $T > 0$ определим
\begin{equation*}
\mathcal{H}_{\text{ode}}(T) = \{ g \circ \varphi_T(\cdot; \theta) : g \in \mathcal{G}, \theta(\cdot) \in L^{\infty}([0, T]; \Theta) \}.
\end{equation*}
Объединяя по всем возможным $T$, получаем основное пространство
\begin{equation*}
\mathcal{H}_{\text{ode}} = \bigcup_{T > 0} \mathcal{H}_{\text{ode}}(T).
\end{equation*}
\end{definition}

Глубине сети $S$ в дискретном случае соответствует время $T$ в непрерывном. Ключевой вопрос: какие функции $F$ можно аппроксимировать с помощью $\mathcal{H}_{\text{ode}}$ (а значит, и с помощью глубоких остаточных сетей (ResNets))?

Важно различать два понятия \cite[Раздел 1.4]{Li2023DeepLD}:
\begin{itemize}
    \item \textbf{Точное представление (Representation):} $F \in \mathcal{H}_{\text{ode}}$.
    \item \textbf{Аппроксимация (Approximation):} $F$ принадлежит замыканию $\mathcal{H}_{\text{ode}}$ в некоторой топологии (например, $L^p$ на компактах). То есть для любого $\varepsilon > 0$ существует $G \in \mathcal{H}_{\text{ode}}$ такая, что $\|F - G\|_{L^p(K)} < \varepsilon$.
\end{itemize}

Потоковые отображения ОДУ обладают сильными топологическими ограничениями: они являются гомеоморфизмами, сохраняющими ориентацию. Следовательно, множество функций, \emph{точно представимых} в виде потока, относительно невелико. Однако это не препятствует тому, чтобы это множество было \emph{плотно} в более широком классе функций (например, в $L^p$), что и означает возможность \emph{аппроксимации}.

Дискретизация Эйлера \eqref{eq:resnet_discrete} аппроксимирует непрерывную динамику \eqref{eq:resnet_ode} с ошибкой $\mathcal{O}(T/S)$ на конечном временном горизонте \cite[Раздел 1.4]{Li2023DeepLD}. Поэтому положительные результаты об аппроксимации для $\mathcal{H}_{\text{ode}}$ автоматически переносятся (с учётом ошибки дискретизации) и на $\mathcal{H}_{\text{resnet}}$.

\subsection{Структурные понятия и основная теорема об универсальной аппроксимации}

В данном подразделе вводится формализм \emph{управляющего} и \emph{терминального семейств}, ключевые структурные понятия \emph{функции-колодца} и \emph{ограниченной аффинной инвариантности}, и формулируется основная теорема об универсальной аппроксимации для размерности $n \geq 2$.


Для компактного описания пространства гипотез $\mathcal{H}_{\text{ode}}$ вводятся следующие определения \cite[Раздел 2.1]{Li2023DeepLD}:

\begin{definition}[Управляющее семейство]
\textbf{Управляющее семейство (control family)} $\mathcal{F}$ --- это множество липшицевых векторных полей $f\colon \mathbb{R}^n \to \mathbb{R}^n$. Каждое поле $f \in \mathcal{F}$ определяет правую часть ОДУ, управляющего динамикой.
\end{definition}

\begin{definition}[Семейство потоковых отображений]
Для горизонта времени $T > 0$ определим семейство потоковых отображений, порождённых $\mathcal{F}$:
\[
\Phi(\mathcal{F}, T) = \{ x \mapsto z(T) : \dot{z}(t) = f_t(z(t)), \ f_t \in \mathcal{F}, \ z(0) = x \},
\]
где $t \mapsto f_t$ --- измеримая и существенно ограниченная функция времени.
\end{definition}

\begin{definition}[Пространство гипотез ODE-сетей]
Для заданных управляющего семейства $\mathcal{F}$ и \textbf{терминального семейства (terminal family)} $\mathcal{G}$ (отображений $\mathbb{R}^n \to \mathcal{Y}$) определим
\[
\mathcal{H}_{\text{ode}}(\mathcal{F}, \mathcal{G}) = \bigcup_{T>0} \{ g \circ \phi : g \in \mathcal{G}, \ \phi \in \Phi(\mathcal{F}, T) \}.
\]
\end{definition}

Таким образом, обучение модели сводится к выбору управляющего поля $f_t$ из $\mathcal{F}$ и финального преобразования $g$ из $\mathcal{G}$.

Для формулировки условий универсальной аппроксимации необходимы два специфических свойства управляющего семейства.

\begin{definition}[Функция-колодец (Well function) \cite{Li2023DeepLD}]
\label{def:well_function}
Липшицева функция $h\colon \mathbb{R}^n \to \mathbb{R}$ называется \textbf{функцией-колодцем}, если существует ограниченное открытое выпуклое множество $U \subset \mathbb{R}^n$ такое, что
\[
\{ x \in \mathbb{R}^n : h(x) = 0 \} = \overline{U}.
\]
Векторнозначная функция $h = (h_1, \ldots, h_{n_0})$ является функцией-колодцем, если каждая её компонента $h_i$ удовлетворяет этому условию.
\end{definition}
\begin{remark}
   Функция $h$ обращается в ноль в точности на компактном выпуклом ``колодце'' $\overline{U}$ и сохраняет постоянный знак вне его. Это позволяет строить динамику, которая ``притягивает'' или ``отталкивает'' траектории относительно этого множества.
\end{remark}

\begin{definition}[Ограниченная аффинная инвариантность (Restricted affine invariance) \cite{Li2023DeepLD}]
\label{def:restricted_affine_invariance}
Управляющее семейство $\mathcal{F}$ называется \textbf{ограниченно аффинно инвариантным}, если для любого $f \in \mathcal{F}$, любого вектора $b \in \mathbb{R}^n$ и любых диагональных матриц $D = \operatorname{diag}(d_1, \ldots, d_n)$ (где $d_i \in \{0, \pm 1\}$) и $A = \operatorname{diag}(a_1, \ldots, a_n)$ (где $|a_i| \leq 1$) выполнено
\[
x \mapsto D f(Ax + b) \in \mathcal{F}.
\]
\end{definition}
Это свойство гарантирует инвариантность семейства относительно изменений знака, покоординатного масштабирования (с коэффициентом $\leq 1$) и сдвигов. Оно выполняется для многих стандартных архитектур, например, для полносвязных слоёв с ReLU-активацией.

Данная теорема \cite[Теорема 2.3]{Li2023DeepLD} даёт общие достаточные условия для того, чтобы пространство гипотез $\mathcal{H}_{\text{ode}}(\mathcal{F}, \mathcal{G})$ обладало свойством универсальной аппроксимации в размерностях $n \geq 2$.

\begin{theorem}[Достаточные условия универсальной аппроксимации, $n \geq 2$]
\label{thm:main_2d}
Пусть $n \geq 2$ и $F\colon \mathbb{R}^n \to \mathbb{R}^m$ --- непрерывная функция. Предположим, что управляющее семейство $\mathcal{F}$ и терминальное семейство $\mathcal{G}$ удовлетворяют условиям:
\begin{enumerate}
    \item \textbf{Покрытие терминальным семейством:} Для любого компакта $K \subset \mathbb{R}^n$ существует липшицева функция $g \in \mathcal{G}$ такая, что $F(K) \subset g(\mathbb{R}^n)$.
    \item \textbf{Наличие функции-колодца в выпуклой оболочке:} Замыкание выпуклой оболочки $\overline{\operatorname{co}}(\mathcal{F})$ содержит функцию-колодец (в смысле Определения \ref{def:well_function}).
    \item \textbf{Ограниченная аффинная инвариантность:} Семейство $\mathcal{F}$ является ограниченно аффинно инвариантным (Определение \ref{def:restricted_affine_invariance}).
\end{enumerate}
Тогда для любого компакта $K \subset \mathbb{R}^n$, любого $p \in [1, \infty)$ и любого $\varepsilon > 0$ существует функция $\hat{F} \in \mathcal{H}_{\text{ode}}(\mathcal{F}, \mathcal{G})$ такая, что
\[
\| F - \hat{F} \|_{L^p(K)} \leq \varepsilon.
\]
\end{theorem}

Для доказательства данной теоремы нам потребуется сформулировать и доказать несколько дополнительных результатов, в том числе одномерную версию теоремы \ref{thm:main_2d}. Для начала разберем идею доказательства.

\begin{proof}[Идея доказательства]
Ключевые шаги доказательства, изложенные в статье \cite[Раздел 4]{Li2023DeepLD}:
\begin{enumerate}
    \item \textbf{Редукция к аппроксимации отображений области (Предложение \ref{prop:reduction} и его следствие):} Если терминальное семейство достаточно богато, то аппроксимацию произвольной $F$ можно свести к аппроксимации некоторого непрерывного отображения $\phi: \mathbb{R}^n \to \mathbb{R}^n$ (преобразования области) потоковыми отображениями с последующим применением $g$.
    \item \textbf{Аппроксимация конечных наборов точек (Леммы \ref{4.13}, \ref{4.14}):} Используя функцию-колодец и ограниченную аффинную инвариантность, можно построить потоковое отображение, которое переводит любой конечный набор различных точек $\{x_k\}$ сколь угодно близко в любой другой набор $\{y_k\}$, сохраняя попарное разделение координат.
    \item \textbf{Переход к $L^p$-аппроксимации (Предложение \ref{prop:high_dim_density}):} Приближая произвольное непрерывное отображение $\phi$ кусочно-постоянным на мелкой сетке и используя лемму о конечных наборах точек для центров ячеек, строим потоковое отображение $\psi$, близкое к $\phi$ в $L^p$.
\end{enumerate}
\end{proof}

\begin{corollary}
    Поскольку явная схема Эйлера аппроксимирует непрерывную динамику с глобальной погрешностью $\mathcal{O}(T/S)$, из плотности $\mathcal{H}_{\text{ode}}$ следует плотность и для соответствующего пространства гипотез дискретных ResNet $\mathcal{H}_{\text{resnet}}$.
\end{corollary}

\subsection{Одномерный случай и его особенности}

Для $n=1$ потоковые отображения обладают специфическим свойством: они являются строго возрастающими функциями \cite[Предложение 3.4]{Li2023DeepLD}.

\begin{proposition}[Монотонность потоковых отображений в 1D.]
\label{prop:1d_monotonicity}
Пусть $n=1$ и $\varphi_T$ --- потоковое отображение автономного ОДУ $\dot{z} = f(z)$ с липшицевой правой частью. Тогда $\varphi_T$ строго возрастает: для любых $x_1 < x_2$ выполнено $\varphi_T(x_1) < \varphi_T(x_2)$.
\end{proposition}
\begin{proof}

Достаточно показать, что если \( z_1 \) и \( z_2 \) --- решения уравнения \ref{eq:resnet_ode}, но с разными начальными значениями \( x_1 < x_2 \), то для всех \( t \geq 0 \) выполняется \( z_1(t) < z_2(t) \).
Предположим противное: \( z_1(t_0) = z_2(t_0) \) при некотором \( t_0 \).
Рассмотрим следующее обыкновенное дифференциальное уравнение:

\[
\dot{w} = -f(w), \quad w(0) = z_1(t_0).
\]

Тогда обе функции \( z_1(t_0 - \cdot) \) и \( z_2(t_0 - \cdot) \) являются решениями этого уравнения. В силу единственности решений получаем
\[
z_1(t_0 - t) = z_2(t_0 - t)
\]
для всех \( t \), откуда следует, что
\[
x_1 = z_1(0) = z_2(0) = x_2,
\]
что противоречит условию \( x_1 < x_2 \). Поскольку обе функции \( z_1 \) и \( z_2 \) непрерывны по \( t \), заключаем, что
\[
z_1(t) < z_2(t)
\]
для всех \( t \).
\end{proof}

\begin{corollary}
    Любая композиция таких отображений также строго возрастает. Следовательно, семейство $\mathcal{A}_{\mathcal{F}} = \bigcup_{T>0} A_{\mathcal{F}}(T)$ (объединение \emph{достижимых множеств}, см. ниже) содержит только возрастающие функции. Это означает, что невозможна аппроксимация функций, убывающих на каком-либо интервале.
\end{corollary}

В рамках класса возрастающих функций универсальная аппроксимация всё же достигается.

\begin{theorem}[Универсальная аппроксимация в 1D \cite{Li2023DeepLD}]
\label{thm:1d_UAT}
Пусть $n = m = 1$, $\mathcal{G} = \{\operatorname{id}\}$ (тождественное отображение). Предположим, что управляющее семейство $\mathcal{F}$ удовлетворяет:
\begin{enumerate}
    \item Оно является ограниченно аффинно инвариантным с $A=1$ (т.е. инвариантно только относительно сдвигов и изменений знака).
    \item $\overline{\operatorname{co}}(\mathcal{F})$ содержит функцию-колодец.
    \item Каждая функция $f \in \mathcal{F}$ является (нестрого) возрастающей.
\end{enumerate}
Тогда для любой непрерывной возрастающей функции $F\colon \mathbb{R} \to \mathbb{R}$ и любого компакта $K \subset \mathbb{R}$ существует последовательность $\hat{F}_k \in \mathcal{H}_{\text{ode}}(\mathcal{F}, \mathcal{G})$ такая, что $\| F - \hat{F}_k \|_{L^p(K)} \to 0$.
\end{theorem}

\textbf{Интерпретация.} В одномерном случае топологические ограничения существенны, и универсальность достигается лишь внутри класса монотонных функций. Это контрастирует с многомерным случаем ($n \geq 2$), где сохранение ориентации не препятствует плотности в $L^p$.

\subsection{Технические основы: свойства потоков и достижимые множества}
Рассмотрим автономное ОДУ 
\[\dot{z(t)} = f(z(t)), \quad z(0)=x,~ f \in \mathcal{F}.\]

\begin{proposition}[Существование, единственность, гладкость \cite{Li2023DeepLD}]
\label{prop:ode_basic}
Если $f$ глобально липшицева, то решение существует, единственно и непрерывно зависит от начального условия. Если $f$ $r$ раз непрерывно дифференцируема ($r \geq 1$), то потоковое отображение $\varphi_T$ является $(r-1)$ раз непрерывно дифференцируемым по начальному условию.
\end{proposition}

\begin{proposition}[Сохраняющие ориентацию диффеоморфизмы \cite{Li2023DeepLD}]
\label{prop:op_diffeomorphism}
Если $f$ дважды непрерывно дифференцируема, то её потоковое отображение $\varphi_T$ является \textbf{сохраняющим ориентацию диффеоморфизмом}, т.е. $\det J_{\varphi_T}(x) > 0$ для всех $x$, где $J_{\varphi_T}$ --- матрица Якоби.
\end{proposition}
Это свойство следует из формулы Лиувилля для якобиана: 
\[J_{\varphi_T}(x) = \exp\left( \int\limits_0^T J_f(\varphi_t(x)) \, dt \right).\]

\begin{proposition}\label{3.6}
Пусть \(n=1\) и \(z(\cdot;x)\) --- решение ОДУ \ref{eq:resnet_ode} с начальным значением \(x\). 
Если \(x\) принадлежит некоторому компактному множеству \(K \subset \mathbb{R}\), 
то конечный модуль непрерывности
\[
\omega_{z(\cdot;x), [0,T]}(r) = 
\sup_{\substack{0 \leq t_1 \leq t_2 \leq T \\ |t_2 - t_1| \leq r}} 
|z(t_1;x) - z(t_2;x)|
\]
стремится к \(0\) при \(r \to 0\) равномерно по \(x \in K\).
\end{proposition}

Для доказательств удобно работать с кусочно-постоянными управлениями.

\begin{definition}[Достижимое множество \cite{Li2023DeepLD}]
\label{def:attainable_set}
Для горизонта $T > 0$ \textbf{достижимое множество} $A_{\mathcal{F}}(T)$ определяется как множество всех композиций конечного числа потоковых отображений, порождённых полями из $\mathcal{F}$:
\[
A_{\mathcal{F}}(T) = \{ \varphi_{t_k}^{f_k} \circ \cdots \circ \varphi_{t_1}^{f_1} : t_1 + \ldots + t_k = T, \ f_1, \ldots, f_k \in \mathcal{F}, \ k \geq 1 \},
\]
где $\varphi_{\tau}^{f}$ --- потоковое отображение за время $\tau$ для поля $f$. Полное \textbf{достижимое множество} есть объединение по всем временам: $A_{\mathcal{F}} = \bigcup_{T>0} A_{\mathcal{F}}(T)$.
\end{definition}
Для аппроксимации достаточно показать, что $A_{\mathcal{F}}$ плотно в пространстве непрерывных отображений.

Важное техническое утверждение показывает, что при построении потоковых отображений можно эффективно использовать выпуклые комбинации полей из $\mathcal{F}$.

\begin{proposition}[Выпуклая релаксация \cite{Li2023DeepLD}]
\label{prop:convex_relaxation}
Пусть $\mathcal{F}$ --- липшицево управляющее семейство, и $\widetilde{\mathcal{F}}$ --- любое липшицево семейство такое, что $\mathcal{F} \subset \widetilde{\mathcal{F}} \subset \overline{\operatorname{co}}(\mathcal{F})$. Тогда достижимые множества совпадают: $A_{\mathcal{F}} = A_{\widetilde{\mathcal{F}}}$.
\end{proposition}
\begin{proof}[Идея доказательства]
Доказательство использует технику быстрого переключения (chattering). Поток, порождённый выпуклой комбинацией полей $\alpha f + (1-\alpha)g$, можно сколь угодно точно приблизить, быстро чередуя эволюцию под действием $f$ и $g$ на малых интервалах времени \cite[Предложение 3.11]{Li2023DeepLD}.
\end{proof}

\begin{remark}
Композиция во времени действует как механизм ``выпуклой релаксации'': возможность составлять сложные траектории из простых позволяет достичь того же эффекта, что и использование более богатого (выпуклого) семейства управлений.
\end{remark}

Предложение \ref{prop:convex_relaxation} является важным результатом, касающимся влияния непрерывной эволюции, которое можно рассматривать как непрерывное семейство композиций: любое семейство функций, управляющее динамической системой, столь же эффективно, как и его выпуклая оболочка, которая может быть значительно более широким семейством функций. Подобные свойства потоков наблюдались в контексте вариационных задач \cite{Warga1962RelaxedVariational}. Это первое указание на пользу композиции в аппроксимации функций.

Для доказательства Предложения \ref{prop:convex_relaxation} потребуются следующие леммы.

\begin{lemma}[Лемма 3.12 \cite{Li2023DeepLD}]\label{lemma:3.12}
Если \(\mathcal{A}_{\mathcal{F}}\), \(\mathcal{A}_{\tilde{\mathcal{F}}}\) --- множества достижимости для \(\mathcal{F}\), \(\tilde{\mathcal{F}}\) и \(\mathcal{F} \subset \tilde{\mathcal{F}} \subset \overline{\mathcal{F}}\), то
\[
\mathcal{A}_{\mathcal{F}} = \mathcal{A}_{\tilde{\mathcal{F}}}.
\]
\end{lemma}

\begin{proof}
Достаточно показать, что \(\mathcal{A}_{\tilde{\mathcal{F}}} \subset \mathcal{A}_{\mathcal{F}}\), тогда будет следовать \(\mathcal{A}_{\mathcal{F}} \subset \mathcal{A}_{\tilde{\mathcal{F}}} \subset \mathcal{A}_{\mathcal{F}}\). Заметим, что любое \(\tilde{\varphi} \in \mathcal{A}_{\tilde{\mathcal{F}}}\) имеет вид
\[
\tilde{\varphi} = \varphi_{t_k}^{\tilde{f}_k} \circ \varphi_{t_{k-1}}^{\tilde{f}_{k-1}} \circ \cdots \circ \varphi_{t_1}^{\tilde{f}_1},
\]
где каждое \(\tilde{f_i}\) принадлежит \(\tilde{\mathcal{F}}\). Чтобы доказать лемму, нужно показать, что для любого компакта \(K \subset \mathbb{R}^n\) и любого \(\varepsilon > 0\) можно построить \(\varphi \in \mathcal{A}_{\mathcal{F}}\) такое, что \(\|\tilde{\varphi} - \varphi\|_{C(K)} \leq \varepsilon\). Докажем это индукцией по \(k \geq 0\). Случай \(k = 0\) очевиден, так как это тождественное отображение. Предположим теперь, что утверждение верно для \(k - 1\), и зафиксируем произвольный компакт \(K\). Запишем \(\tilde{\varphi} = \varphi_{t_k}^{\tilde{f}_k} \circ \tilde{\psi}\), где \(\tilde{\psi}\) --- композиция \(k - 1\) потоковых отображений, управляемых семейством \(\tilde{\mathcal{F}}\). По индукционному предположению, для любого \(\varepsilon_1\) (которое будет выбрано позже) существует \(\psi \in \mathcal{A}_{\mathcal{F}}\) такое, что \(\|\psi - \tilde{\psi}\|_{C(K)} \leq \varepsilon_1\). Более того, в силу предположения \(\tilde{\mathcal{F}} \subset \mathcal{F}\), для любого \(\varepsilon_2\) и компакта \(K'\) существует функция \(f_k \in \mathcal{F}\) такая, что \(\|f_k - \tilde{f_k}\|_{C(K')} \leq \varepsilon_2\). Здесь мы выбираем \(K' = \{x : \inf_{\psi \in \mathcal{V}(K)} \|x - y\| \leq 2(\varepsilon_1 + t_k)e^{t_k \cdot \mathrm{Lip}(\tilde{f_k})}\}\) и \(\varepsilon_2 < 1\).

Теперь рассмотрим два обыкновенных дифференциальных уравнения:
\[
\tilde{z}(t) = \tilde{\psi}(x) + \int\limits_0^t f_k(\tilde{z}(\tau)) \, d\tau \quad\text{и}\quad z(t) = \psi(x) + \int\limits_0^t f_k(z(\tau)) \, d\tau.
\]
Имеем \(\tilde{\varphi}(x) = \tilde{z}(t_k)\) и отображение \(x \mapsto z(t_k)\) принадлежит \(\mathcal{A}_{\mathcal{F}}\), следовательно, остаётся показать, что решения этих ОДУ можно сделать сколь угодно близкими. Фактически, покажем, что
\[
\sup_{0 \leq t \leq t_k} |z(t) - \tilde{z}(t)| < 2(\varepsilon_1 + t_k \varepsilon_2)e^{t_k \cdot \mathrm{Lip}(\tilde{f_k})}. \tag{$\ast$}
\]
Предположим противное, тогда положим
\[
t = \inf\{s \geq 0 : |z(s) - \tilde{z}(s)| \geq 2(\varepsilon_1 + t_k \varepsilon_2)e^{t_k \cdot \mathrm{Lip}(\tilde{f_k})}\}.
\]

По непрерывности имеем \(|z(t) - \bar{z}(t)| \geq 2(\varepsilon_1 + t_k \varepsilon_2) e^{t_k \mathrm{Lip}(\bar{f}_k)}\). По предположению также имеем \(|\bar{f}_k(z(\tau)) - f_k(z(\tau))| \leq \varepsilon_2\) для всех \(\tau < t\). Вычитая, получаем
\[
|\bar{z}(t) - z(t)| \leq \varepsilon_1 + \left| \int\limits_{0}^{t} \bar{f}_k(\bar{z}(\tau)) d\tau - \int\limits_{0}^{t} f_k(z(\tau)) d\tau \right|
\]
\[
\leq \varepsilon_1 + \int\limits_{0}^{t} |\bar{f}_k(z(\tau)) - f_k(z(\tau))| d\tau + \int\limits_{0}^{t} |\bar{f}_k(\bar{z}(\tau)) - \bar{f}_k(z(\tau))| d\tau
\]
\[
\leq \varepsilon_1 + \varepsilon_2 t + \mathrm{Lip}(\bar{f}_k) \int\limits_{0}^{t} |\bar{z}(\tau) - z(\tau)| d\tau.
\]

По неравенству Гронуолла имеем
\[
\sup_{0 \leq t \leq t_k} |\bar{z}(t) - z(t)| \leq (\varepsilon_1 + t_k \varepsilon_2) e^{t_k \mathrm{Lip}(\bar{f}_k)},
\]
что противоречит выбору \(t\). Следовательно, ($\ast$) выполнено, и мы можем выбрать \(\varepsilon_1\) и \(\varepsilon_2\) сколь угодно малыми, что завершает доказательство.
\end{proof}

\begin{lemma}[Лемма 3.13 \cite{Li2023DeepLD}]\label{lemma:3.13}
Предположим, что \(f, g \in \mathcal{F}\) и \(t > 0\). Тогда \(\varphi_t^{(f+g)/2} \in \mathcal{A}_{\mathcal{F}}\).
\end{lemma}

\begin{proof}
Покажем, что \(\varphi_{t/2N}^f \circ \varphi_{t/2N}^g \circ \cdots \circ \varphi_{t/2N}^f \circ \varphi_{t/2N}^g\) может сколь угодно хорошо приближать \(\varphi_t^{(f+g)/2} \in \mathcal{A}_{\mathcal{F}}\) при увеличении \(N\). Отображение \(\varphi_t^{(f+g)/2}\) является решением
\begin{multline*}
    z(t) = x + \int\limits_{0}^{t} \left( \frac{f + g}{2} \right) (z(\tau)) d\tau = x + \left( \int\limits_{0}^{t/2N} + \int\limits_{2t/2N}^{3t/2N} + \cdots + \int\limits_{(2N-2)t/2N}^{(2N-1)t/2N} \right) (f + g)(z(\tau)) d\tau + \\
    + \left( \int\limits_{t/2N}^{2t/2N} + \cdots + \int\limits_{(2N-1)t/2N}^{t} \right) \left[ \left( \frac{f + g}{2} \right) (z(\tau)) - \left( \frac{f + g}{2} \right) \left( z \left( \tau - \frac{t}{2N} \right) \right) \right] d\tau = \\
    = x + \left( \int\limits_{0}^{t/2N} + \int\limits_{2t/2N}^{3t/2N} + \cdots + \int\limits_{(2N-2)t/2N}^{(2N-1)t/2N} \right) f(z(\tau)) d\tau + \left( \int\limits_{t/2N}^{2t/2N} + \cdots + \int\limits_{(2N-1)t/2N}^{t}\right) g(z(\tau)) d\tau + \\
    + \left( \int\limits_{t/2N}^{2t/2N} + \cdots + \int\limits_{(2N-1)t/2N}^{t} \right) \left[ \left( \frac{f + g}{2} \right) (z(\tau)) - \left( \frac{f + g}{2} \right) \left( z \left( \tau - \frac{t}{2N} \right) \right) \right] + \\
    + \left[ g \left( z \left( \tau - \frac{t}{2N} \right) \right) - g(z(\tau)) \right] d\tau.
\end{multline*}

Таким образом, если
\begin{multline*}
    w(t) = x + \left( \int\limits_{0}^{t/2N} + \int\limits_{2t/2N}^{3t/2N} + \cdots + \int\limits_{(2N-2)t/2N}^{(2N-1)t/2N} \right) f(w(\tau)) d\tau +\\
    + \left( \int\limits_{t/2N}^{2t/2N} + \cdots + \int\limits_{(2N-1)t/2N}^{t}\right) g(w(\tau)) d\tau,
\end{multline*}
то
\begin{multline*}
    |z(t) - w(t)| \leq \int\limits_{0}^{t} \max(\mathrm{Lip}(f), \mathrm{Lip}(g))|z(\tau) - w(\tau)| \, d\tau + \\ 
    + \frac{t}{2}\omega_{z,[0,t]} \left( \frac{t}{2N} \right) \left[ \mathrm{Lip} \left( \frac{f + g}{2} \right) + \mathrm{Lip}(g) \right].
\end{multline*}

Напомним, что \(\omega\) --- модуль непрерывности, определённый в Предложении \ref{3.6}. Снова применяя неравенство Гронуолла, получаем
\[
|z(t) - w(t)| \leq \frac{t}{2}\omega_{z,[0,t]} \left( \frac{t}{2N} \right) \left[ \mathrm{Lip} \left( \frac{f + g}{2} \right) + \mathrm{Lip}(g) \right] e^{\max(\mathrm{Lip}(f), \mathrm{Lip}(g))}.
\]

Для любого выбранного компакта \(K\), \(\omega_{z,[0,t]} (\frac{t}{2n}) \to 0\) по Предложению \ref{3.6}, следовательно, получаем \(\varphi_{t}^{(f+g)/2} \in \mathcal{A}_{\mathcal{F}}\).
\end{proof}

Теперь мы готовы доказать Предложение \ref{prop:convex_relaxation}.

\begin{proof}[Доказательство Предложения \ref{prop:convex_relaxation}]
Используя ту же технику, что и в доказательстве Леммы \ref{lemma:3.13}, можно показать, что для \(f_1, \ldots, f_m \in \mathcal{F}\) имеем \(\varphi_{t}^h \in \mathcal{A}_{\mathcal{F}}\), где \(h = \sum_i q_i f_i\) для некоторых рациональных чисел \(q_i\). Пусть \(\mathcal{A}_{\mathcal{F}}\) --- множество достижимости с управляющим семейством \(\mathcal{F}' = \{\sum_{i=1}^{m} q_i f_i : q_i \in \mathbb{Q}, \sum_i q_i = 1, f_i \in \mathcal{F}, m \in \mathbb{N}\}\). Тогда \(\mathcal{A}_{\mathcal{F}} = \mathcal{A}_{\mathcal{F}}\). Поскольку \(\mathcal{F}' = \mathcal{F}\), мы приходим к желаемому результату, используя Лемму \ref{lemma:3.12}.
\end{proof}

\subsubsection{Редукция к аппроксимации преобразований области (Слайд 15)}
\label{subsubsec:reduction_to_domain}

Следующий результат крайне важен, так как позволяет свести задачу аппроксимации произвольной функции $F$ к задаче аппроксимации некоторого отображения $\phi: \mathbb{R}^n \to \mathbb{R}^n$ (преобразования области) потоковыми отображениями.

\begin{proposition}[Редукция (3.8) \cite{Li2023DeepLD}]
\label{prop:reduction}
Пусть $F: \mathbb{R}^n \to \mathbb{R}^m$ непрерывна, $K \subset \mathbb{R}^n$ компакт, и существует липшицева $g \in \mathcal{G}$ такая, что $F(K) \subset g(\mathbb{R}^n)$. Тогда для любого $\varepsilon > 0$ существует непрерывное отображение $\phi: K \to \mathbb{R}^n$ такое, что $\| F - g \circ \phi \|_{L^p(K)} \leq \varepsilon$.
\end{proposition}

\begin{proof}
Следует из общего результата о композиции функций, доказанного в \cite{28}. 
Мы доказываем этот специальный случай здесь для полноты изложения.

Множество \( F(K) \) компактно, поэтому для любого \( \delta > 0 \) мы можем образовать разбиение 
\( F(K) = \bigcup_{i=1}^{N} B_i \) с диаметром \( \operatorname{diam}(B_i) \leq \delta \). 
По предположению, \( g^{-1}(B_i) \) непусто для каждого \( i \), поэтому выберем 
\( z_i \in g^{-1}(B_i) \). Для каждого \( i \) определим \( A_i = F^{-1}(B_i) \cap K \), 
так что \(\{A_i\}\) образует разбиение \( K \). 
По внутренней регулярности меры Лебега, для любого \( \delta' > 0 \) и для каждого \( i \) 
можно найти компактное множество \( K_i \subset A_i \) такое, что 
\( \lambda(A_i \setminus K_i) \leq \delta' \) (где \( \lambda \) --- мера Лебега) 
и множества \( K_i \) попарно не пересекаются. 
По лемме Урысона, для каждого \( i \) существует непрерывная функция 
\( \varphi_i : K \to [0,1] \) такая, что \( \varphi_i = 1 \) на \( K_i \) и 
\( \varphi_i = 0 \) на \( \bigcup_{j \neq i} K_j \). 
Теперь образуем непрерывную функцию

\[
\varphi(x) = \sum_{i=1}^{N} z_i \varphi_i(x).
\]

Определим \( K' = \left\{\sum_{i=1}^{N} \alpha_i z_i : \alpha_i \in [0,1]\right\} \), 
которое компактно и \( \varphi(K) \subset K' \). Тогда

\begin{align*}
\|F - g \circ \varphi\|_{L^p(K)} 
&\leq \sum_{i=1}^{N} \|F - g \circ \varphi\|_{L^p(K_i)} 
   + \sum_{i=1}^{N} \|F - g \circ \varphi\|_{L^p(A_i \setminus K_i)} \\
&\leq \sum_{i=1}^{N} \|F - g \circ \varphi_i\|_{L^p(K_i)} 
   + \|F\|_{C(K)} + \|g\|_{C(K')} N \delta'.
\end{align*}

Мы выбираем \( \delta' \) достаточно малым, чтобы последнее слагаемое было ограничено величиной \( \delta \). Тогда

\[
\|F - g \circ \varphi\|_{L^p(K)} \leq \sum_{i=1}^{N} \delta \lambda(K_i) + \delta 
\leq (1 + \lambda(K)) \delta.
\]

Взяв \( \delta = \varepsilon / (1 + \lambda(K)) \), получаем требуемый результат.
\end{proof}

\begin{corollary}
\label{cor:reduction_corollary}
Если семейство потоковых отображений $\Phi(\mathcal{F}, T)$ плотно в множестве непрерывных отображений $\phi: \mathbb{R}^n \to \mathbb{R}^n$ (в топологии $L^p$ на компактах), то и $\mathcal{H}_{\text{ode}}(\mathcal{F}, \mathcal{G})$ плотно в множестве непрерывных $F$.
\end{corollary}
Таким образом, основная техническая задача сводится к аппроксимации \emph{произвольного} непрерывного $\phi$ отображениями из $A_{\mathcal{F}}$.

\begin{remark}
В одномерном пространстве $A_{\mathcal{F}}$ содержит только возрастающие функции.    
\end{remark}

Наша цель --- показать, что при выполнении условий Теоремы \ref{thm:1d_UAT} можно аппроксимировать \emph{любую} непрерывную возрастающую функцию.
Стратегия:
\begin{enumerate}
    \item Построить динамику перестановок (\emph{rearrangement dynamics}), способную точно перевести любой конечный набор точек $\{x_i\}$ в любой другой набор $\{y_i\}$ (при условии $x_1 < \ldots < x_M$, $y_1 < \ldots < y_M$).
    \item Использовать плотность кусочно-линейных интерполянтов для аппроксимации произвольной непрерывной возрастающей функции.
\end{enumerate}

\begin{example}[Cемейство ReLU и функция-колодец]
    Рассмотрим управляющее семейство для одномерной полносвязной ResNet с ReLU-активацией:
    \[
    \mathcal{F}_{\text{ReLU}} = \{ x \mapsto v \cdot \operatorname{ReLU}(wx + b) : v, w, b \in \mathbb{R} \}, \quad \operatorname{ReLU}(z) = \max(0, z).
    \]
    Оно, очевидно, ограниченно аффинно инвариантно. Явный пример функции-колодца с нулевым множеством $[q_1, q_2]$:
    \[
    h_Q(x) = \frac{1}{2} \left[ \operatorname{ReLU}(q_1 - x) + \operatorname{ReLU}(x - q_2) \right] \in \overline{\operatorname{co}}(\mathcal{F}_{\text{ReLU}}).
    \]
    Это демонстрирует выполнение абстрактного условия ``$\overline{\operatorname{co}}(\mathcal{F})$ содержит функцию-колодец'' на конкретном примере.    
\end{example}

Ключевая лемма для одномерного случая.
\begin{lemma}[Конечная интерполяция в 1D \cite{Li2023DeepLD}]
\label{lem:finite_interp_1d}
Пусть $\mathcal{F}$ удовлетворяет условиям Теоремы \ref{thm:1d_UAT}. Тогда для любых возрастающих последовательностей $x_1 < \ldots < x_M$ и $y_1 < \ldots < y_M$ существует потоковое отображение $\psi \in A_{\mathcal{F}}$ такое, что $\psi(x_i) = y_i$ для всех $i=1,\ldots,M$.
\end{lemma}
\begin{proof}[Идея доказательсвта]
Индукция по $M$. Используя функцию-колодец $h_Q$ (с нулём на $[q_1, q_2]$), можно построить динамику, которая:
\begin{enumerate}
    \item С помощью сдвигов и отражений (используя аффинную инвариантность) располагает точку $x_1$ и цель $y_1$ по одну сторону от ``колодца''.
    \item Применяя поток поля $\pm h_Q$, перемещает $x_1$ в $y_1$, не затрагивая другие точки (т.к. они находятся в нулевом множестве $h_Q$ или с другой стороны).
    \item Индуктивно повторяет процесс для оставшихся точек.
\end{enumerate}
\end{proof}

\begin{proposition}
\label{prop:1d_density}
Пусть $\phi: \mathbb{R} \to \mathbb{R}$ непрерывна и возрастает, $K=[a,b]$ --- компакт. Тогда для любого $\varepsilon > 0$ существует $\psi \in A_{\mathcal{F}}$ такая, что $\| \phi - \psi \|_{C(K)} < \varepsilon$.
\end{proposition}
\begin{proof}[Идея доказательства]
\leavevmode
\begin{enumerate}
    \item Выбираем достаточно мелкое разбиение $a = x_1 < x_2 < \ldots < x_M = b$ так, чтобы колебание $\phi$ на каждом подинтервале не превышало $\varepsilon$ (используем модуль непрерывности $\omega_\phi$).
    \item Полагаем $y_i = \phi(x_i)$. По Лемме \ref{lem:finite_interp_1d} строим $\psi \in A_{\mathcal{F}}$, интерполирующую эти точки: $\psi(x_i) = y_i$.
    \item В силу монотонности $\phi$ и $\psi$, на каждом интервале $[x_i, x_{i+1}]$ выполняется $|\phi(x) - \psi(x)| \leq \omega_\phi(|x_{i+1} - x_i|) \leq \varepsilon$.
\end{enumerate}
\end{proof}
Эта предложение вместе со следствием \ref{cor:reduction_corollary} доказывает Теорему \ref{thm:1d_UAT}.

Займёмся вопросами скорости аппроксимации в одномерном случае.
Для семейства ReLU можно получить количественные оценки сложности аппроксимации в терминах времени горизонта $T$.
\begin{definition}[Сложность через полную вариацию]
Для гладкой возрастающей функции $\phi: [0,1] \to \mathbb{R}$ определим
\[
T_0(\phi) := \| \ln \phi' \|_{\mathrm{TV}[0,1]},
\]
где $\|\cdot\|_{\mathrm{TV}}$ --- полная вариация, а $\phi'$ --- производная.
\end{definition}
\begin{proposition}[Связь времени горизонта и сложности \cite{Li2023DeepLD}]
\label{prop:1d_complexity}
Пусть $\phi$ --- кусочно-гладкая возрастающая функция на $[0,1]$.
\begin{enumerate}
    \item Если $T \geq T_0(\phi)$, то $\phi$ может быть \emph{точно представлена} как потоковое отображение с временным горизонтом $\leq T$.
    \item Если $T < T_0(\phi)$, то существует нижняя оценка ошибки аппроксимации, зависящая от разницы $T_0(\phi) - T$.
\end{enumerate}
\end{proposition}
\textit{Интерпретация}: величина $\| \ln \phi' \|_{\mathrm{TV}}$ измеряет ``объем изменений'', необходимый для представления функции $\phi$ потоком. Она играет роль, аналогичную числу слоёв в дискретной сети.

\subsection{Многомерный случай ($n \geq 2$): преодоление ограничений}
В размерности $n \geq 2$ потоковые отображения по-прежнему являются сохраняющими ориентацию гомеоморфизмами. Однако ключевое отличие от одномерного случая состоит в следующем известном факте:
\begin{quote}
Множество сохраняющих ориентацию диффеоморфизмов \textbf{плотно} в $L^p$ на ограниченных областях при $n \geq 2$ \cite{Brenier2003LpApproximation}.
\end{quote}
Таким образом, топологическое ограничение на точное представление не мешает плотности в смысле аппроксимации. Задача сводится к тому, чтобы показать, что достижимое множество $A_{\mathcal{F}}$ достаточно богато, чтобы приближать произвольные непрерывные отображения.

Основная техническая трудность в многомерном случае --- построить потоковое отображение, которое одновременно управляет многими точками. Решение основывается на двух леммах.

\begin{lemma}[Разделение координат \cite{Li2023DeepLD}]
\label{lem:coordinate_separation}
Предположим, $\mathcal{F}$ содержит функцию-колодец. Тогда для любых различных точек $x_1, \ldots, x_m \in \mathbb{R}^n$ и $\varepsilon > 0$ существует потоковое отображение $\psi \in A_{\mathcal{F}}$ такое, что:
\begin{enumerate}
    \item $|\psi(x_k) - x_k| < \varepsilon$ для всех $k$,
    \item Для каждой координаты $i$ значения $\{ [\psi(x_k)]_i \}_{k=1}^m$ являются $m$ различными вещественными числами.
\end{enumerate}
\end{lemma}
\textit{Идея}: используя функцию-колодец и аффинные преобразования, можно построить динамику, которая слегка ``сдвигает'' точки так, чтобы их координаты стали попарно различными, не смешивая при этом изначально различные координаты.

\begin{lemma}[Точечное управление \cite{Li2023DeepLD}]
\label{lem:pointwise_steering}
Пусть точки $x_1, \ldots, x_m$ удовлетворяют заключению Леммы \ref{lem:coordinate_separation}. Тогда для любых целевых точек $y_1, \ldots, y_m \in \mathbb{R}^n$ и $\varepsilon > 0$ существует $\phi \in A_{\mathcal{F}}$ такое, что $|\phi(x_k) - y_k| < \varepsilon$ для всех $k$.
\end{lemma}
\textit{Идея}: работая покоординатно и используя тот факт, что координаты точек разделены, можно независимо перемещать каждую точку вдоль выбранной координатной оси с помощью подходящим образом сконструированных из функции-колодца полей.

Композиция отображений, построенных в этих леммах, позволяет аппроксимировать произвольное соответствие на конечном наборе точек.

Финальный шаг --- переход от аппроксимации на конечных множествах к аппроксимации в $L^p$ \cite[Предложение 4.11]{Li2023DeepLD}.

\begin{proposition}[Универсальная аппроксимация в $n \geq 2$]
\label{prop:high_dim_density}
Пусть $n \geq 2$, $\mathcal{F}$ ограниченно аффинно инвариантно и $\overline{\operatorname{co}}(\mathcal{F})$ содержит функцию-колодец. Тогда для любого компакта $K \subset \mathbb{R}^n$, любого $p \in [1, \infty)$, любого непрерывного $\phi: \mathbb{R}^n \to \mathbb{R}^n$ и любого $\varepsilon > 0$ существует $\psi \in A_{\mathcal{F}}$ такое, что
\[\| \phi - \psi \|_{L^p(K)} < \varepsilon.\]
\end{proposition}

\begin{proof}[Идея доказательства]
\begin{enumerate}
    \item Аппроксимируем $\phi$ на $K$ кусочно-постоянным отображением $\hat{\phi}$ на мелкой кубической сетке.
    \item Сопоставим каждому кубу сетки его центр $x_k$ и значение $\hat{\phi}(x_k) = y_k$.
    \item Используя Леммы \ref{lem:coordinate_separation} и \ref{lem:pointwise_steering}, построим потоковое отображение $\psi$, переводящее каждое $x_k$ близко в $y_k$.
    \item Показываем, что из-за малости размера сетки и непрерывности $\phi$ отображение $\psi$ будет близко к $\phi$ на всём $K$ в норме $L^p$.
\end{enumerate}
\end{proof}

Заметим, что для целей аппроксимации тот факт, что \(\overline{\text{co}}(\mathcal{F})\) содержит функцию типа ``колодца'' \(h\), позволяет нам без потери общности считать, что \(\mathcal{F}\) содержит такую функцию. Чтобы убедиться в этом, обозначим через \(\tilde{\mathcal{F}}\) наименьшее ограниченно аффинно инвариантное множество, содержащее \(\mathcal{F} \cup \{h\}\). Мы имеем \(\mathcal{F} \subset \tilde{\mathcal{F}} \subset \overline{\text{co}}(\mathcal{F})\). Предложение 3.11 тогда утверждает, что \(\mathcal{A}_{\mathcal{F}} = \mathcal{A}_{\tilde{\mathcal{F}}}\), следовательно, мы можем доказывать результаты аппроксимации, используя \(\tilde{\mathcal{F}}\) вместо \(\mathcal{F}\).

Для доказательства Предложения \ref{prop:high_dim_density} нам потребуется несколько предварительных результатов. Ключевой подход аналогичен одномерному случаю: мы показываем, что можем преобразовать конечное число различных исходных точек в конечное число целевых точек, которые не обязательно различны. Начнем со следующей леммы, которая обобщает Лемму \ref{lem:finite_interp_1d}, приведённую в одномерном случае.

\begin{lemma}[4.12]\label{4.12}
Предположим, что \(\mathcal{F}\) содержит функцию типа ``колодца''. Пусть \(\varepsilon > 0\) и \(x^1, \ldots, x^m, y^1, \ldots, y^m \in \mathbb{R}^n\) таковы, что \(\{x^k\}\) являются различными точками. Тогда существует \(\psi \in \mathcal{A}_{\mathcal{F}}\) такой, что \(|\psi(x^k) - y^k| \leq \varepsilon\) для всех \(k = 1, \ldots, m\).
\end{lemma}

Лемма \ref{4.12} следует из комбинации следующих двух лемм.

\begin{lemma}[4.13]\label{4.13}
Предположим, что \(\mathcal{F}\) содержит функцию типа ``колодца'' и \(x^1, \ldots, x^m\) --- различные точки. Тогда для любого \(\varepsilon > 0\) существует поток \(\psi \in \mathcal{A}_{\mathcal{F}}\) такой, что \(|\psi(x^k) - x^k| \leq \varepsilon\) и для каждого \(i = 1, \ldots, n\), \([\psi(x^k)]_i\) (i-я координата \(\psi(x^k)\)), \(k = 1, \ldots, m\), представляют собой \(m\) различных вещественных чисел.
\end{lemma}

\begin{lemma}[4.14]\label{4.14}
Предположим, что \(\mathcal{F}\) содержит функцию типа ``колодца'', \(x^1, \ldots, x^m\) --- различные точки и удовлетворяют результату Леммы \ref{4.13}, то есть \(\{x_i^k\}\) являются \(m\) различными вещественными числами для любого \(i\). Тогда для любого \(\varepsilon > 0\) и для \(m\) целевых точек \(y^1, \ldots, y^m\) существует поток \(\psi \in \mathcal{A}_{\mathcal{F}}\) такой, что \(|\psi(x^k) - y^k| \leq \varepsilon\).
\end{lemma}

Поскольку \( K \) компактно, по теореме о продолжении достаточно рассмотреть случай, когда \( K \) является гиперкубом. Мы можем для простоты взять единичный гиперкуб \( K = [0, 1]^n \), так как общий случай аналогичен. Для \( \varphi \in L^p(K) \), согласно стандартной теории аппроксимации, \( \varphi \) может быть аппроксимирована кусочно постоянными функциями, т.е. существует
\[
\hat{\varphi} = \sum_i \varphi_i \chi_i 
\]
такая, что \(\|\hat{\varphi} - \varphi\|_{L^p(K)} \leq \varepsilon/2\). Здесь \( i = [i_1, \ldots, i_n] \) --- мультииндекс, \( \varphi_i \in \mathbb{R}^n \) и \( \chi_i \) --- индикатор куба
\[
\square_i = \left[ \frac{i_1}{N}, \frac{i_1 + 1}{N} \right] \times \cdots \times \left[ \frac{i_n}{N}, \frac{i_n + 1}{N} \right]. 
\]
Мы также обозначим \( p_i = (i_1/N, \ldots, i_n/N) \). Определим сжатый куб
\[
\square_i^\alpha = \left[ \frac{i_1}{N}, \frac{i_1 + \alpha}{N} \right] \times \cdots \times \left[ \frac{i_n}{N}, \frac{i_n + \alpha}{N} \right], 
\]
где \( 0 < \alpha \leq 1 \). Мы имеем \( K = \bigcup_i \square_i \), и определим \( K^\alpha = \bigcup_i \square_i^\alpha \). Мы также построим одномерную функцию сжатия, \( h^\alpha : [0, 1] \rightarrow [0, 1] \), такую что \( h^\alpha(x) = i/N \) если \( i/N \leq x \leq (i + \alpha)/N \), и непрерывно возрастающую на \([0, 1]\). Используя это, мы можем сформировать \( n \)-мерное отображение сжатия через тензорное произведение:
\[
H^\alpha(x) = (h^\alpha(x_1), \ldots, h^\alpha(x_n)). 
\]

Идея доказательства Предложения \ref{prop:high_dim_density} довольно проста: мы просто сжимаем каждую ячейку сетки \( \square_i \) приблизительно в точку \( p_i \), затем используем лемму выше для преобразования каждой \( p_i \) в \( \varphi_i \). Последнее обсуждается в предварительном шаге; мы строим "почти" сжимающее отображение в \( \mathcal{A}_{\mathcal{F}} \), которое аппроксимирует \( H^\alpha \).

\textbf{Утверждение.} Для заданной точности \( \varepsilon_1 > 0 \) существует поток \( \tilde{H} \in \mathcal{A}_{\mathcal{F}} \) такой, что
\[
\| \tilde{H} - H^\alpha \|_{C(K)} \leq \varepsilon_1.
\]

\begin{proof}[Доказательство утверждения]
Поскольку \( h^\alpha \) возрастает и непрерывна, мы хотим использовать наш результат в одномерном случае. Конкретно, мы демонстрируем, как ограничить \( n \)-мерное семейство управлений до одномерного.

Предположим, что \(\mathcal{F}\) является \( n \)-мерным семейством управлений. Тогда для каждого \( f = (f_1, \ldots, f_n) \in \mathcal{F} \) мы определим динамику, управляемую его ограничением на первую координату, как
\[
\dot{z}_1 = f_1(x_1), \quad \dot{z}_i = 0 \text{ для } i \geq 2, 
\]
т.е. возьмём \( D = A = \text{diag}(1, 0, \ldots, 0) \). Множество таких систем управления обозначается через \( \mathcal{F}_{R,1} \) (\( R \) означает ограничение, а 1 означает первую координату). Очевидно, что \( \mathcal{F}_{R,1} \) замкнуто относительно композиции. Более того, \( \mathcal{A}_{\mathcal{F}_{R,1}} \) совпадает с
\[
\mathcal{A}_{\mathcal{R}} \times \{\text{id}\} \times \{\text{id}\} \times \cdots \times \{\text{id}\}, 
\]
где \(\mathcal{R}\) --- одномерное семейство управлений
\[
\mathcal{R} = \{g(x) : g(x) = f_1(x, 0, \ldots, 0), \, f \in \mathcal{F}\}. 
\]

Поскольку \(\mathcal{F}\) содержит функцию типа ``колодца'', то и \(\mathcal{R}\) содержит такую функцию. По Предложению \ref{prop:1d_density} мы можем найти \(\tilde{h} \in \mathcal{R}\) такую, что \(\| \tilde{h} - h^{\alpha} \|_{C([0,1])} \leq \varepsilon / n\).

По композиции мы знаем, что \(\tilde{H} := (\tilde{h}, \ldots, \tilde{h})\) принадлежит \(\mathcal{A}_{\mathcal{F}}\), и \(\|\tilde{H} - H^{\alpha}\|_{C(K)} \leq \varepsilon\).
\end{proof}

Мы используем вышеуказанные обозначения: \(\psi\) для переноса \(p_i\) в \(\varphi_i\), и \(\tilde{H}\) для приближённого сжимающего отображения, удовлетворяющих следующим оценкам:
\[
|\psi(p_i) - \varphi_i| \leq \varepsilon_1 < 1, 
\]
\[
\|\tilde{H} - H^{\alpha}\|_{C(K)} \leq \varepsilon_2 < 1,
\]
где \(\varepsilon_1\) и \(\varepsilon_2\) будут определены позже.

Теперь оценим погрешность \(\|\psi \circ \tilde{H} - \varphi\|_{L^{p}(K)}\). Для любого \(\alpha\) мы можем записать
\begin{align*}
\|\psi \circ \tilde{H} - \varphi\|_{L^{p}(K)} &\leq \|\psi \circ \tilde{H} - \hat{\varphi}\|_{L^{p}(K)} + \|\hat{\varphi} - \varphi\|_{L^{p}(K)} \\
&\leq \|\psi \circ \tilde{H} - \hat{\varphi}\|_{L^{p}(K^{\alpha})} + \|\psi \circ \tilde{H} - \hat{\varphi}\|_{L^{p}(K \setminus K^{\alpha})} + \varepsilon/2 \\
&\leq J_1 + J_2 + \varepsilon/2. 
\end{align*}

\textbf{Оценка \(J_1\).} Имеем
\begin{align*}
J_1 &= \|\psi \circ \tilde{H} - \psi \circ H^{\alpha}\|_{L^{p}(K^{\alpha})} + \|\psi \circ H^{\alpha} - \hat{\varphi}\|_{L^{p}(K^{\alpha})} \\
&\leq \omega_{\psi}(\|\tilde{H} - H^{\alpha}\|_{C(K)}) + \sum_{i} |\psi(p_i) - \varphi_i| \cdot |\square_i|^{1/p} \\
&\leq \omega_{\psi}(\|\tilde{H} - H^{\alpha}\|_{C(K)}) + N^{n-n/p} \varepsilon_1 \\
&\leq \omega_{\psi}(\varepsilon_2) + N^{n-n/p} \varepsilon_1.
\end{align*}

\textbf{Оценка \(J_2\).} Обозначим через \(\tilde{K} = [-1, 2]^n\) расширенный куб. Имеем
\[
J_2 \leq |K \setminus K^{\alpha}| \cdot (\text{diam}(\psi(\tilde{K})) + \|\hat{\varphi}\|_{L^{\infty}(K)}). 
\]

Мы выбираем \(\psi\) так, что \(\varepsilon_1 \leq \varepsilon/8\), \(\alpha\) так, что
\[
|K \setminus K^{\alpha}| \leq (\text{diam}(\psi(\tilde{K})) + \|\hat{\varphi}\|_{L^{\infty}(K)})^{-1} \varepsilon/4, 
\]
и, наконец, \(\tilde{H}\) так, что \(\omega_{\psi}(\varepsilon_2) \leq \varepsilon/8\). Тогда \(\|\psi \circ \tilde{H} - \varphi\|_{L^{p}(K)} \leq \varepsilon\), и мы полагаем \(\tilde{\varphi} = \psi \circ \tilde{H} \in \mathcal{A}_{\mathcal{F}}\), что даёт требуемый результат.

Как и в одномерном случае, Предложение \ref{prop:high_dim_density} вместе со следствием предложения \ref{prop:reduction} доказывают Теорему \ref{thm:main_2d}.

\subsection{Семейства тензорных произведений и сравнение с классической теорией}

Иногда удобно строить многомерные управления, ``складывая'' одномерные.

\begin{definition}[Семейство тензорного произведения]
Пусть $\mathcal{F}^{(1)}$ --- одномерное управляющее семейство. Определим $n$-мерное \textbf{семейство тензорного произведения}:
\[
\widetilde{\mathcal{F}} = \{ x \mapsto (g(x_1), g(x_2), \ldots, g(x_n)) : g \in \mathcal{F}^{(1)} \}.
\]
\end{definition}
Если затем замкнуть $\widetilde{\mathcal{F}}$ относительно всех аффинных преобразований (а не только диагональных), получим более богатое семейство $\mathcal{F}^{(n)}$.

\begin{proposition}[Универсальность тензорных произведений \cite{Li2023DeepLD}]
\label{prop:tensor_universality}
Предположим, что $A_{\mathcal{F}^{(1)}}$ содержит все непрерывные возрастающие функции на $\mathbb{R}$. Тогда при $n \geq 2$ соответствующее семейство $\mathcal{F}^{(n)}$ обеспечивает универсальную $L^p$-аппроксимацию на $\mathbb{R}^n$.
\end{proposition}

Подход динамических систем предлагает новую парадигму аппроксимации.
\begin{itemize}
    \item \textbf{Классическая линейная/нелинейная аппроксимация:} Используется набор $D$ простых функций (полиномы, вейвлеты). Аппроксимант --- \textbf{линейная комбинация} конечного числа элементов словаря. Сложность определяется числом слагаемых $N$ и затуханием/разреженностью коэффициентов \cite{DeVore1998NonlinearApproximation, Daubechies1992TenLectures, Mallat1998WaveletTour}.
    \item \textbf{Подход динамических систем:} Сложность возникает в основном за счёт \textbf{глубины композиции} (времени горизонта $T$), а не линейных комбинаций. Управляющее семейство $\mathcal{F}$ может быть очень маленьким (например, содержать одну функцию-колодец в своей выпуклой оболочке). Потоковые отображения --- это композиции малых возмущений тождества, управляемых $\mathcal{F}$.
\end{itemize}
Таким образом, мы имеем новую парадигму: ``\emph{аппроксимация через композицию потоков}'' вместо ``аппроксимации через линейную комбинацию статических базисных функций''.

\subsection{Следствия для глубоких сетей и связь с теорией управления}
Проведённый анализ проясняет, когда ResNet-подобные архитектуры являются универсальными аппроксиматорами.
\begin{itemize}
    \item \textbf{Структурные условия носят архитектурный характер:}
    \begin{itemize}
        \item Наличие в управляющем семействе возможности формировать \emph{функции-колодцы} (что выполняется для ReLU, сигмоиды, $\tanh$ и др.).
        \item Инвариантность относительно базовых аффинных преобразований (знака, сдвига, масштабирования), что обычно встроено в полносвязные и свёрточные слои.
    \end{itemize}
    \item \textbf{Многие практические архитектуры попадают в эту схему:}
    \begin{itemize}
        \item Полносвязные ResNet с активациями ReLU, sigmoid, $\tanh$.
        \item Остаточные блоки, состоящие из нескольких слоёв (residual blocks).
        \item Подходящие тензорно-произведные конструкции.
    \end{itemize}
    \item \textbf{Источник выразительной силы --- композиция:} каждый элементарный блок может быть очень простым, а сложное поведение возникает за счёт их последовательной композиции во времени (глубине).
\end{itemize}

Рассматриваемая проблема имеет глубокие связи с другими областями математики.
\begin{itemize}
    \item \textbf{Теория управления:} Аппроксимация с помощью $A_{\mathcal{F}}$ связана с задачей \textbf{управляемости (controllability)} \cite{Sussmann1990NonlinearControllability, Balachandran2002ControllabilitySurvey, Chukwu1991ControllabilityAbstract}. Однако классическая управляемость фокусируется на переводе \emph{одной} точки в целевую, тогда как здесь требуется управлять \emph{целым множеством} входных значений $K$ с помощью одного и того же управления $f_t$. Это можно рассматривать как задачу управляемости в бесконечномерном пространстве функций.
    \item \textbf{Динамические системы:} Результаты дополняют классические работы о том, когда гомеоморфизм \emph{точно представим} как потоковое отображение \cite{Fort1955EmbeddingHomeomorphisms, Utz1981EmbeddingHomeomorphisms, Zhang2009EmbeddingDiffeomorphisms}, и о плотности сохраняющих ориентацию диффеоморфизмов в $L^p$ \cite{Brenier2003LpApproximation}. Наши результаты дают конструктивный способ аппроксимации произвольных отображений потоками с ограниченной структурой управления.
    \item \textbf{Оптимальное управление:} Обучение в этой формулировке можно переформулировать как задачу \textbf{управления в среднем поле (mean-field optimal control)} \cite{E2019MeanFieldOptimalControl, Li2017MaximumPrinciple}, что открывает путь к применению принципа максимума Понтрягина и методов Гамильтона-Якоби-Беллмана.
\end{itemize}

Подведём итоги лекции:
\begin{itemize}
    \item Глубокие остаточные сети допускают естественную интерпретацию как управляемые потоки ОДУ.
    \item При выполнении мягких структурных условий (наличие функции-колодца в выпуклой оболочке управляющего семейства, ограниченная аффинная инвариантность и простое терминальное семейство) соответствующее пространство гипотез $\mathcal{H}_{\text{ode}}$ является универсальным аппроксиматором в $L^p$ для размерностей $n \geq 2$.
    \item В размерности $n=1$ универсальность достигается в классе непрерывных возрастающих функций, причём можно получить оценки скорости аппроксимации через полную вариацию $\| \ln \phi' \|_{\mathrm{TV}}$.
    \item Эти результаты подчёркивают выразительную силу \textbf{композиции через потоки} и устанавливают связь между глубоким обучением, теорией динамических систем и оптимальным управлением.
\end{itemize}
Для дальнейшего ознакомления предлагается обратиться к оригиналам статей, упомянутых в конспекте.
\end{document}